\documentclass[10pt,a4paper]{amsart}
\usepackage[margin=19mm]{geometry}
\usepackage{amsmath,amssymb,mathtools}
\usepackage[scr=rsfs,cal=euler]{mathalfa}
\usepackage{xfrac}
\usepackage[unicode,hidelinks,bookmarks=false]{hyperref}
\newcommand{\ie}{i.e.\ }
\newcommand{\cf}{cf.\ }
\newcommand{\resp}{respectively\ }
\newcommand{\Subsection}{\subsection}
\providecommand{\nobreakdash}{\nobreak\hbox{-}}
\setlength{\emergencystretch}{2em}
\multlinegap0pt
\usepackage{amssymb}
\usepackage{mathtools}
\usepackage{enumerate}
\usepackage[normalem]{ulem}
\usepackage{bbm}
\newtheorem{lemma}{Lemma}
\title{Exceptional zeros of $\mathrm{GL}_3\times\mathrm{GL}_3$ Rankin--Selberg $L$-functions}
\author{Jesse Thorner}
\date{Source: arXiv:2508.09984v2 (CC BY 4.0)}
\begin{document}
\maketitle

\noindent\emph{Source and licence.} Exceptional zeros of $\mathrm{GL}_3\times\mathrm{GL}_3$ Rankin--Selberg $L$-functions. Version arXiv:2508.09984v2 (CC BY 4.0), distributed by arXiv under the Creative Commons Attribution 4.0 International licence, http://creativecommons.org/licenses/by/4.0/, as recorded in the arXiv licence metadata for this exact version and shown on the abstract page https://arxiv.org/abs/2508.09984v2. Full licence notice: \texttt{licenses/arxiv-2508.09984-CC-BY-4.0.txt}.\par\smallskip

\noindent\textit{Editorial scope.} Counted unit: the article's lemma lem:coeff (source0.tex lines 385--388). The article's complete original proof of it is delivered with it (source0.tex lines 389--417, one \texttt{proof} environment of 29 lines). No mathematics is written, repaired or re-derived: the statement and the proof are the article's own lines, in the article's own notation, and no retained line is substituted or re-flowed.  Retained uncounted as the source-local context the proof needs: lines 217--226; lines 290--305.  Nothing was omitted from any retained span, so the package carries no omission record.  No retained line contains a citation, so the wrapper carries no bibliography and no bibliography entry.  No retained line contains a figure environment, an image-inclusion command or drawing code, so no image is delivered and the package declares no asset dependency.  Editorial formatting only, in addition: an AMS \texttt{amsart} preamble that restates the article's own declarations and macros used by the retained lines, namely the environments \texttt{lemma} on separate counters, together with the standard packages those lines need.  The retained environments print in this document as Lemma 1 in order of appearance, whereas the article prints the same items with the numbering of its own full document.  The complete native source of the article as distributed in the arXiv e-print for this exact version is bundled unchanged as \texttt{sources/183207/source0.tex}.
	\begin{equation}
	\label{eqn:a_def}
	\begin{aligned}
	a_{\pi\times\pi'}(v^{\ell})&= \begin{cases}
 	a_{\pi}(v^{\ell})a_{\pi'}(v^{\ell})&\mbox{if $v\notin S_{\pi}\cup S_{\pi'}$,}\\
 	\sum_{j=1}^n \sum_{j'=1}^{n'}\alpha_{j,j',\pi\times\pi'}(v)^{\ell}&\mbox{if $v\in S_{\pi}\cup S_{\pi'}$,}
 \end{cases}\\
 a_{\widetilde{\pi}\times\widetilde{\pi}'}(v^{\ell})&=\overline{a_{\pi\times\pi'}(v^{\ell})}.
 \end{aligned}
	\end{equation}

\begin{lemma}
\label{lem:CG}
Let $j,k,u\geq 0$ be integers, $\pi\in\mathfrak{F}_{2}$, and $\psi_1,\psi_2,\chi\in\mathfrak{F}_{1}$.  Suppose that if
\[
u\in\{j,k\}\cup\{j+k-2r\colon 0\leq r\leq \min\{j,k\}\},
\]
then $\mathrm{Sym}^u(\pi)\in\mathfrak{F}_{u+1}$.  Then
\[
\mathrm{Sym}^{j}(\pi\otimes\psi_1)\boxtimes\mathrm{Sym}^k(\pi\otimes\psi_2)\otimes\chi=\mathop{\boxplus}_{r=0}^{\min\{j,k\}}\mathrm{Sym}^{j+k-2r}(\pi)\otimes\chi\psi_1^j \psi_2^k\omega_{\pi}^{r}\in\mathfrak{A}_{(j+1)(k+1)},
\]
and $L(s,\mathrm{Sym}^{j}(\pi\otimes\psi_1)\boxtimes(\mathrm{Sym}^k(\pi\otimes\psi_2)\otimes\chi))=L(s,\mathrm{Sym}^{j}(\pi\otimes\psi_2)\times(\mathrm{Sym}^k(\pi\otimes\psi_2)\otimes\chi))$. In particular, if $v\nmid\infty$, then the following identities hold:
\begin{align}
a_{\mathrm{Sym}^{j}(\pi)\times(\mathrm{Sym}^k(\pi)\otimes\chi)}(v^{\ell}) &= \sum_{r=0}^{\min\{j,k\}}a_{\mathrm{Sym}^{j+k-2r}(\pi)\otimes\chi \omega_{\pi}^{r}}(v^{\ell}),\label{eqn:CG1}\\
a_{\mathrm{Sym}^{j}(\pi)\times\mathrm{Sym}^j(\widetilde{\pi})}(v^{\ell})&=1+\sum_{r=1}^j a_{\mathrm{Sym}^{2r}(\pi)\otimes\omega_{\pi}^{-r}}(v^{\ell}).\label{eqn:CG2}
\end{align}
\end{lemma}

\begin{lemma}
	\label{lem:coeff}
Let $(\pi,\pi',\chi)\in\mathfrak{F}_2\times\mathfrak{F}_2\times\mathfrak{F}_1$ with $\pi\not\sim\pi'$.  If $v\notin S_{\pi}\cup S_{\pi'}\cup S_{\chi}$, then $a_{\mathcal{D}}(v^{\ell})\geq 0$.
\end{lemma}

\begin{proof}
Fix $\ell\geq 1$ and a place $v\nmid\infty$ of $F$ such that $v\notin S_{\pi}\cup S_{\pi'}\cup S_{\chi}$ We use \eqref{eqn:a_def} throughout the proof without further mention.  Since $v$ and $\ell$ are fixed, we suppress $(v^{\ell})$ throughout (e.g., $a_{\mathcal{D}}=a_{\mathcal{D}}(v^{\ell})$ and $a_{\mathrm{Ad}(\pi)}=a_{\mathrm{Ad}(\pi)}(v^{\ell})$).  Since $\mathrm{Ad}(\pi)$ and $\mathrm{Ad}(\pi')$ are self-dual, it follows that $a_{\mathrm{Ad}(\pi)},a_{\mathrm{Ad}(\pi')}\in\mathbb{R}$.  A direct computation shows that $a_{\mathcal{D}}$ equals
\begin{align*}
&6+4a_{\mathrm{Ad}(\pi)\times(\mathrm{Ad}(\pi')\otimes\chi)}+4a_{\mathrm{Ad}(\pi)\times(\mathrm{Ad}(\pi')\otimes\overline{\chi})}+7a_{\mathrm{Ad}(\pi)}+2a_{\mathrm{Ad}(\pi')}\\
&+2a_{\mathrm{Ad}(\pi')\otimes\chi}+2a_{\mathrm{Ad}(\pi')\otimes\overline{\chi}}+5a_{\mathrm{Sym}^4(\pi)\otimes\overline{\omega}_{\pi}^{2}} + 2a_{\mathrm{Sym}^4(\pi')\otimes\overline{\omega}_{\pi'}^{2}} + 3a_{\mathrm{Ad}(\pi)\times\mathrm{Ad}(\pi')} \\
&+ 3a_{\mathrm{Ad}(\pi)\times(\mathrm{Sym}^4(\pi')\otimes\overline{\omega}_{\pi'}^{2})}+ a_{\mathrm{Ad}(\pi')\times(\mathrm{Sym}^4(\pi)\otimes\overline{\omega}_{\pi}^{2})}+ 2a_{\mathrm{Ad}(\pi')\times(\mathrm{Sym}^4(\pi)\otimes\chi \overline{\omega}_{\pi}^{2})}\\
&+2a_{\mathrm{Ad}(\pi')\times(\mathrm{Sym}^4(\pi)\otimes\overline{\chi}\,\overline{\omega}_{\pi}^{2})} + a_{\mathrm{Sym}^4(\pi)\times(\mathrm{Sym}^4(\pi')\otimes\overline{\omega}_{\pi}^{2} \overline{\omega}_{\pi'}^{2})}.
\end{align*}
Since $v\notin S_{\pi}\cup S_{\pi'}\cup S_{\chi}$, we can rewrite our expression for $a_{\mathcal{D}}$ as
\begin{align*}
&2 a_{\mathrm{Ad}(\pi')\otimes\chi} (a_{\mathrm{Sym}^4(\pi)\otimes\overline{\omega}_{\pi}^2}+a_{\mathrm{Ad}(\pi)}+1)+2 a_{\mathrm{Ad}(\pi')\otimes\overline{\chi}} (a_{\mathrm{Sym}^4(\pi)\otimes\overline{\omega}_{\pi}^2}+a_{\mathrm{Ad}(\pi)}+1)\\
&+2 a_{\mathrm{Ad}(\pi)} a_{\mathrm{Ad}(\pi')\otimes\chi}+2 a_{\mathrm{Ad}(\pi)} a_{\mathrm{Ad}(\pi')\otimes\overline{\chi}}+2 a_{\mathrm{Ad}(\pi)} (a_{\mathrm{Sym}^4(\pi')\otimes\overline{\omega}_{\pi'}^2}\\
&+a_{\mathrm{Ad}(\pi')}+1)+4 (a_{\mathrm{Sym}^4(\pi)\otimes\overline{\omega}_{\pi}^2}+a_{\mathrm{Ad}(\pi)}+1)+(a_{\mathrm{Sym}^4(\pi')\otimes\overline{\omega}_{\pi'}^2}+a_{\mathrm{Ad}(\pi')}+1)\\
&+(a_{\mathrm{Sym}^4(\pi)\otimes\overline{\omega}_{\pi}^2}+a_{\mathrm{Ad}(\pi)}+1)(a_{\mathrm{Sym}^4(\pi')\otimes\overline{\omega}_{\pi'}^2}+a_{\mathrm{Ad}(\pi')}+1)
\end{align*}
It follows from Lemma \ref{lem:CG} that if $v\notin S_{\pi}\cup S_{\pi'}\cup S_{\chi}$, then
\begin{align*}
a_{\mathrm{Ad}(\pi)}^2=1+a_{\mathrm{Ad}(\pi)}+a_{\mathrm{Sym}^4(\pi)\otimes\overline{\omega}_{\pi}^2},\qquad a_{\mathrm{Ad}(\pi')}^2=1+a_{\mathrm{Ad}(\pi')}+a_{\mathrm{Sym}^4(\pi')\otimes\overline{\omega}_{\pi'}^2}.
\end{align*}
Therefore, our expression for $a_{\mathcal{D}}$ simplifies to
\begin{align*}
&2a_{\mathrm{Ad}(\pi)}^2 a_{\mathrm{Ad}(\pi')\otimes\chi}+2a_{\mathrm{Ad}(\pi)}^2 a_{\mathrm{Ad}(\pi')\otimes\overline{\chi}}+2a_{\mathrm{Ad}(\pi)}a_{\mathrm{Ad}(\pi')\otimes\chi}+2a_{\mathrm{Ad}(\pi)}a_{\mathrm{Ad}(\pi')\otimes\overline{\chi}}\\
&\qquad +2 a_{\mathrm{Ad}(\pi)} a_{\mathrm{Ad}(\pi')}^2+4 a_{\mathrm{Ad}(\pi)}^2+a_{\mathrm{Ad}(\pi')}^2+a_{\mathrm{Ad}(\pi)}^2 a_{\mathrm{Ad}(\pi')}^2\\
&=4 \mathrm{Re}(a_{\text{Ad}(\pi)}^2 a_{\text{Ad}(\pi')\otimes\chi}+a_{\text{Ad}(\pi)} a_{\text{Ad}(\pi')\otimes\chi})\\
&\qquad+2 a_{\text{Ad}(\pi)} |a_{\text{Ad}(\pi')\otimes\chi}|^2+4 a_{\text{Ad}(\pi)}^2+|a_{\text{Ad}(\pi')\otimes\chi}|^2+a_{\text{Ad}(\pi)}^2 |a_{\text{Ad}(\pi')\otimes\chi}|^2\\
&=|2a_{\mathrm{Ad}(\pi)}+a_{\mathrm{Ad}(\pi)}a_{\mathrm{Ad}(\pi')\otimes\chi}+a_{\mathrm{Ad}(\pi')\otimes\chi}|^2,
\end{align*}
which is nonnegative.
\end{proof}

\end{document}
